% Solutions of 24-graph-traversals.
\section*{Chapter \ref{ch:b1:graph-traversals} --- Graphs and Traversals}

\begin{solution}{exo:b1:graph-traversals:1}
$\deg 1 = 2$, $\deg 2 = 3$, $\deg 3 = 3$, $\deg 4 = 2$, $\deg 5 = 3$, $\deg 6 = 1$: the sum is $14 = 2 \times 7$. The
\omterm{def:b1:graph-traversals:graph}{vertices} of odd \omterm{def:b1:graph-traversals:degree}{degree} are 2, 3, 5 and 6: four of them, an even number.
\end{solution}

\begin{solution}{exo:b1:graph-traversals:2}
Times $d|f$: $1$: $1|12$, $2$: $2|5$, $4$: $3|4$, $3$: $6|11$, $5$: $7|10$, $6$: $8|9$. \omterm{def:b1:graph-traversals:edge-classes}{Tree edges} $1 \to 2$, $2 \to 4$,
$1 \to 3$, $3 \to 5$, $5 \to 6$; \omterm{def:b1:graph-traversals:edge-classes}{back edges} $4 \to 1$ and $6 \to 5$; \omterm{def:b1:graph-traversals:edge-classes}{cross edge} $3 \to 4$ (4 is finished and was
discovered before 3); no \omterm{def:b1:graph-traversals:edge-classes}{forward edge}.
\end{solution}

\begin{solution}{exo:b1:graph-traversals:3}
$\delta(1, 1) = 0$; $\delta(1, 2) = \delta(1, 3) = 1$; $\delta(1, 4) = \delta(1, 5) = 2$; $\delta(1, 6) = 3$. With
neighbours taken in increasing order, the \omterm{def:b1:graph-traversals:bfs}{BFS tree} has the \omterm{def:b1:graph-traversals:graph}{edges} $12, 13, 24, 35, 56$.
\end{solution}

\begin{solution}{exo:b1:graph-traversals:4}
With the tasks in the order they first appear, Kahn's method gives boil, chop, tin, pasta, fry, sauce, serve. A
\omterm{def:b1:graph-traversals:dfs}{depth-first search} gives, by decreasing \omterm{def:b1:graph-traversals:dfs}{finishing time}, tin, chop, fry, sauce, boil, pasta, serve. Both put each task
after the ones it needs.
\end{solution}

\begin{solution}{exo:b1:graph-traversals:5}
A \omterm{def:b1:graph-traversals:graph}{vertex} is unmarked while it is open. On a DAG the search never meets an open \omterm{def:b1:graph-traversals:graph}{vertex} again (that would be a \omterm{def:b1:graph-traversals:edge-classes}{back
edge}), and a finished \omterm{def:b1:graph-traversals:graph}{vertex} is marked: each \omterm{def:b1:graph-traversals:graph}{vertex} is visited once, 31 visits for 31 \omterm{def:b1:graph-traversals:graph}{vertices}. On a \omterm{def:b1:graph-traversals:path}{cycle}, the search
comes back to an open \omterm{def:b1:graph-traversals:graph}{vertex}, which is not marked, and starts it again: $x$, $y$, $x$, $y$, \dots, until the call \omterm{def:b1:c-functions-and-stack:stack-overflow}{stack
overflows}. Marking on entry, as \code{dfs} does, fixes it.
\end{solution}

\begin{solution}{exo:b1:graph-traversals:6}
The search starts from the first \omterm{def:b1:graph-traversals:graph}{vertex} only, so it colours the \omterm{def:b1:graph-traversals:path}{path} 1--2--3--4 and never sees the triangle, an odd
\omterm{def:b1:graph-traversals:path}{cycle}. Starting a search from every uncoloured \omterm{def:b1:graph-traversals:graph}{vertex}, as \code{two_colour} does, repairs it; the cost stays
$O(n + m)$.
\end{solution}

\begin{solution}{exo:b1:graph-traversals:7}
\omcode{code/bachelor-1/24-graph-traversals/exercises/count_paths.py}{6}{17}{Counting shortest \omterm{def:b1:graph-traversals:path}{paths}}
A shortest \omterm{def:b1:graph-traversals:path}{path} to $v$ ends with an edge $u v$ where $\delta(s, u) = \delta(s, v) - 1$, and those \omterm{def:b1:graph-traversals:path}{paths} are counted by
the sum of the counts of such $u$. Every such $u$ leaves the queue before $v$ does, and each of its predecessors left
before it: its count is final when it is added. The test compares with an enumeration of the \omterm{def:b1:graph-traversals:path}{paths} on random \omterm{def:b1:graph-traversals:graph}{graphs}.
\end{solution}

\begin{solution}{exo:b1:graph-traversals:8}
\omcode{code/bachelor-1/24-graph-traversals/exercises/find_cycle.py}{4}{28}{A \omterm{def:b1:graph-traversals:path}{cycle}, from the first \omterm{def:b1:graph-traversals:edge-classes}{back edge}}
By \cref{prop:b1:graph-traversals:back}, a \omterm{def:b1:graph-traversals:path}{cycle} exists exactly when the search meets a \omterm{def:b1:graph-traversals:edge-classes}{back edge} $u \to v$; the parents
lead from $u$ up to $v$, and reversed they give the \omterm{def:b1:graph-traversals:path}{cycle}. The search is $O(n + m)$, and the \omterm{def:b1:graph-traversals:path}{cycle} is built once.
\end{solution}

\begin{solution}{exo:b1:graph-traversals:9}
By induction on $n$; true for $n = 1$. A \omterm{def:b1:graph-traversals:free-tree}{free tree} with $n \geq 2$ \omterm{def:b1:graph-traversals:graph}{vertices} has a \omterm{def:b1:graph-traversals:graph}{vertex} of \omterm{def:b1:graph-traversals:degree}{degree} 1: the end $x$ of a
longest \omterm{def:b1:graph-traversals:path}{simple path} has a neighbour on the \omterm{def:b1:graph-traversals:path}{path}, and no other, since another neighbour on the \omterm{def:b1:graph-traversals:path}{path} would close a \omterm{def:b1:graph-traversals:path}{cycle}
and one off it would give a longer \omterm{def:b1:graph-traversals:path}{path}. Removing $x$ and its edge leaves a \omterm{def:b1:graph-traversals:free-tree}{free tree} on $n - 1$ \omterm{def:b1:graph-traversals:graph}{vertices}, with $n - 2$
\omterm{def:b1:graph-traversals:graph}{edges} by induction: $n - 1$ in all. Conversely, take a \omterm{def:b1:graph-traversals:connected}{connected graph} with $n$ \omterm{def:b1:graph-traversals:graph}{vertices} and $n - 1$ \omterm{def:b1:graph-traversals:graph}{edges}; while it
has a \omterm{def:b1:graph-traversals:path}{cycle}, remove an edge of a \omterm{def:b1:graph-traversals:path}{cycle}, which keeps it \omterm{def:b1:graph-traversals:connected}{connected}. The result is a \omterm{def:b1:graph-traversals:free-tree}{free tree}, with $n - 1$ \omterm{def:b1:graph-traversals:graph}{edges}: no
edge was removed, so there was no \omterm{def:b1:graph-traversals:path}{cycle}.
\end{solution}

\begin{solution}{exo:b1:graph-traversals:10}
We show by induction that each search of the second pass assigns exactly one \omterm{def:b1:graph-traversals:connected}{strongly connected component}. Suppose the
assigned \omterm{def:b1:graph-traversals:graph}{vertices} are a union of components, and the search starts from $s$, the unassigned \omterm{def:b1:graph-traversals:graph}{vertex} with the largest
\omterm{def:b1:graph-traversals:dfs}{finishing time}, in the component $C$. Reversing \omterm{def:b1:graph-traversals:graph}{arcs} keeps components, and none of $C$ is assigned, so the search
reaches all of $C$. Suppose it leaves $C$: the first \omterm{def:b1:graph-traversals:graph}{vertex} reached outside $C$, in a component $C'$ not assigned, is
reached by a reversed \omterm{def:b1:graph-traversals:graph}{arc}, that is, an \omterm{def:b1:graph-traversals:graph}{arc} of $G$ from $C'$ to $C$. By \cref{lem:b1:graph-traversals:finish} some
\omterm{def:b1:graph-traversals:graph}{vertex} of $C'$ finishes after every \omterm{def:b1:graph-traversals:graph}{vertex} of $C$, so after $s$: impossible, since $s$ has the largest \omterm{def:b1:graph-traversals:dfs}{finishing time}
among the unassigned \omterm{def:b1:graph-traversals:graph}{vertices}. So the search assigns exactly $C$, and the induction goes on.
\end{solution}

\begin{solution}{exo:b1:graph-traversals:11}
\omcode{code/bachelor-1/24-graph-traversals/exercises/euler.py}{6}{23}{Hierholzer's method}
Since every \omterm{def:b1:graph-traversals:degree}{degree} is even, a walk that enters a \omterm{def:b1:graph-traversals:graph}{vertex} other than its start can leave it again: the walk gets stuck
only at the \omterm{def:b1:graph-traversals:graph}{vertex} where its current part began. When the top of the stack has no edge left, it is placed in the
circuit, and the walk resumes from the \omterm{def:b1:graph-traversals:graph}{vertex} below; the circuit, read backwards, is a single closed walk with every
detour spliced in where it started. Each edge is pushed and popped once: $O(n + m)$. The test builds \omterm{def:b1:graph-traversals:graph}{graphs} as unions
of random closed walks, whose \omterm{def:b1:graph-traversals:degree}{degrees} are all even, and checks that the circuit uses each edge once.
\end{solution}

\begin{solution}{exo:b1:graph-traversals:12}
If $G$ is two-coloured, the colours alternate along a \omterm{def:b1:graph-traversals:path}{cycle}, which returns to its first colour after an even number of
steps. Conversely, suppose $G$ has no odd \omterm{def:b1:graph-traversals:path}{cycle}, and colour each component by the parity of the \omterm{def:b1:graph-traversals:distance}{distance} from a \omterm{def:b1:graph-traversals:graph}{vertex}
$s$ found by \omterm{def:b1:graph-traversals:bfs}{breadth-first search}. If an edge $u v$ joins two \omterm{def:b1:graph-traversals:graph}{vertices} of the same colour, let $w$ be the last common
\omterm{def:b1:graph-traversals:graph}{vertex} of the tree \omterm{def:b1:graph-traversals:path}{paths} from $s$ to $u$ and to $v$. The tree \omterm{def:b1:graph-traversals:path}{paths} from $w$ to $u$ and to $v$ share only $w$, and their
lengths $\delta(s, u) - \delta(s, w)$ and $\delta(s, v) - \delta(s, w)$ have the same parity; with $u v$ they form a
\omterm{def:b1:graph-traversals:path}{cycle} of odd length, a contradiction.
\end{solution}

\begin{solution}{pb:b1:graph-traversals:1}
\textbf{1.} None: the smallest \omterm{def:b1:graph-traversals:degree}{degree} is 0, and person 40 has it. The average \omterm{def:b1:graph-traversals:degree}{degree} is $100 / 40 = 2.5$.

\textbf{2.} A component is a class of the relation ``reachable from''; nobody is reachable from person 40 but
person 40.

\textbf{3.} The groups 1--14 and 28--39 are joined by two friendships ($1$--$28$ and $5$--$37$) into one component,
together with three people of the middle group (18, 20 and 25), whose friends are on that side. The ten others of the
middle group form a component of their own.

\textbf{4.} A triangle is a \omterm{def:b1:graph-traversals:path}{cycle} of length 3; the colours would alternate along it and return to the first colour
after three steps, which is impossible.

\textbf{5.} Yes: a \omterm{def:b1:graph-traversals:graph}{graph} is \omterm{def:b1:graph-traversals:bipartite}{bipartite} exactly when all its components are. The program would print an odd \omterm{def:b1:graph-traversals:path}{cycle} of the
component that is not.

\textbf{6.} $O(n + m)$: each \omterm{def:b1:graph-traversals:graph}{vertex} enters the queue or the stack once, and each list of neighbours is read once.

\textbf{7.} 106: two for each of the 53 roads.

\textbf{8.} Every \omterm{def:b1:graph-traversals:graph}{arc} but $1 \to 3$ comes with its reverse, so every place still reaches the places it reached before,
by going around road 1 if needed: from the Library, by the South \omterm{def:b1:graph-traversals:bridge}{Bridge} back to the Town Square.

\textbf{9.} Along any \omterm{def:b1:graph-traversals:path}{path} the place numbers increase, so no \omterm{def:b1:graph-traversals:path}{path} returns to its start: there is no \omterm{def:b1:graph-traversals:path}{cycle}. Two
distinct \omterm{def:b1:graph-traversals:graph}{vertices} that reach each other would make a \omterm{def:b1:graph-traversals:path}{cycle}, so each component has one \omterm{def:b1:graph-traversals:graph}{vertex}.

\textbf{10.} The \omterm{def:b1:graph-traversals:condensation}{component graph} is the DAG itself; the \omterm{def:b1:graph-traversals:graph}{vertices} without incoming \omterm{def:b1:graph-traversals:graph}{arcs} are the places never listed
second: 1, 2, 4, 5, 6, 8 and 18.

\textbf{11.} $O(n + m)$ with $n = 30$ and $m = 105$, for each of the two searches and for reversing the \omterm{def:b1:graph-traversals:graph}{arcs}.

\textbf{12.} A \omterm{def:b1:graph-traversals:path}{path} from $u$ to $v$ in $G$ is, reversed, a \omterm{def:b1:graph-traversals:path}{path} from $v$ to $u$ in the reversed \omterm{def:b1:graph-traversals:graph}{graph}; reaching each
other is unchanged.

\textbf{13.} If $u v$ lies on a \omterm{def:b1:graph-traversals:path}{cycle}, the rest of the \omterm{def:b1:graph-traversals:path}{cycle} joins $u$ and $v$ after removing it, and every \omterm{def:b1:graph-traversals:path}{path} that
used the edge can go around: the components stay the same. If $u v$ lies on no \omterm{def:b1:graph-traversals:path}{cycle}, any other \omterm{def:b1:graph-traversals:path}{path} from $u$ to $v$
would close a \omterm{def:b1:graph-traversals:path}{cycle} with it; so removing it separates $u$ from $v$.

\textbf{14.} When an edge $u v$ is first examined from $u$, $v$ is either unmarked, and the edge becomes a \omterm{def:b1:graph-traversals:edge-classes}{tree edge}, or
open, an ancestor; if $v$ were finished, $u v$ would already have been examined from $v$, when $u$ was unmarked or
open, so $u$ would be a descendant of $v$.

\textbf{15.} With no \omterm{def:b1:graph-traversals:edge-classes}{cross edges}, every edge leaving the \omterm{def:b1:induction-and-termination:rooted-tree}{subtree} of $v$ other than $u v$ goes to an ancestor of $v$, so
to $u$ or above, with a \omterm{def:b1:graph-traversals:dfs}{discovery time} at most $d(u)$. Such an edge exists exactly when $\mathrm{low}(v) \leq d(u)$,
and it closes a \omterm{def:b1:graph-traversals:path}{cycle} through $u v$ with tree \omterm{def:b1:graph-traversals:path}{paths}. If none exists, every \omterm{def:b1:graph-traversals:path}{path} from the \omterm{def:b1:induction-and-termination:rooted-tree}{subtree} to the rest uses
$u v$, which is a \omterm{def:b1:graph-traversals:bridge}{bridge}.

\textbf{16.} The \omterm{def:b1:graph-traversals:edge-classes}{tree edge} itself would give $\mathrm{low}(v) \leq d(u)$ and hide every \omterm{def:b1:graph-traversals:bridge}{bridge}. Only once, because a
second edge between the same two \omterm{def:b1:graph-traversals:graph}{vertices} is a real \omterm{def:b1:graph-traversals:path}{cycle} of two \omterm{def:b1:graph-traversals:graph}{edges}: in a road map, two parallel roads.

\textbf{17.} Ten of them join a person with a single friend, who would be cut off by losing that friendship; $12$--$20$
holds people 20 and 18 to the rest, and $33$--$39$ people 33 and 35. Inside the groups, no single friendship matters.

\textbf{18.} Joining each place to its three nearest makes \omterm{def:b1:graph-traversals:path}{cycles} likely, but two clusters could still be joined by a
single road: it is not a proof. The program checked: 0 \omterm{def:b1:graph-traversals:bridge}{bridges}.

\textbf{19.} Let $u v$ be a \omterm{def:b1:graph-traversals:bridge}{bridge}, oriented $u \to v$. Every \omterm{def:b1:graph-traversals:path}{path} from the side of $v$ to the side of $u$ must use the
edge, which points the other way: no orientation makes $u$ reachable from $v$.

\textbf{20.} An edge is examined from both ends. A \omterm{def:b1:graph-traversals:edge-classes}{tree edge} is added when the parent examines it; from the child, it
is skipped, since the \omterm{def:b1:graph-traversals:graph}{arc} already exists. A non-tree edge joins a \omterm{def:b1:graph-traversals:graph}{vertex} to an ancestor; it is added when the
descendant examines it, the ancestor being discovered earlier, and skipped from the ancestor, which has the larger
\omterm{def:b1:graph-traversals:dfs}{discovery time} on that side of the test, or the \omterm{def:b1:graph-traversals:graph}{arc} already exists.

\textbf{21.} \omterm{def:b1:graph-traversals:edge-classes}{Tree edges} point from parent to child, so the tree \omterm{def:b1:graph-traversals:path}{path} from the root to any \omterm{def:b1:graph-traversals:graph}{vertex} is a directed \omterm{def:b1:graph-traversals:path}{path}.

\textbf{22.} Let $v$ be a \omterm{def:b1:graph-traversals:graph}{vertex} other than the root, with parent $p$. The edge $p v$ is not a \omterm{def:b1:graph-traversals:bridge}{bridge}, so by question
15 some edge goes from a \omterm{def:b1:graph-traversals:graph}{vertex} $x$ of the \omterm{def:b1:induction-and-termination:rooted-tree}{subtree} of $v$ to a proper ancestor $a$ of $v$. It is oriented upwards,
$x \to a$, and tree \omterm{def:b1:graph-traversals:graph}{arcs} lead from $v$ down to $x$: $a$ is reachable from $v$, and $d(a) < d(v)$. Repeating from $a$,
the \omterm{def:b1:graph-traversals:dfs}{discovery times} decrease, so the walk ends at the root.

\textbf{23.} A \omterm{def:b1:graph-traversals:bridge}{bridge} forbids a \omterm{def:b1:graph-traversals:connected}{strongly connected} orientation (question 19). Without a \omterm{def:b1:graph-traversals:bridge}{bridge}, the orientation of
\code{orient} lets every \omterm{def:b1:graph-traversals:graph}{vertex} reach the root and the root reach every \omterm{def:b1:graph-traversals:graph}{vertex}, so every \omterm{def:b1:graph-traversals:graph}{vertex} reaches every other.

\textbf{24.} The search starts at the Town Square, place 1, and its first neighbour is the Library, place 3, since
road 1 comes first in the table: $1 \to 3$ is a \omterm{def:b1:graph-traversals:edge-classes}{tree edge}, oriented from parent to child.

\textbf{25.} \emph{Robbins' theorem}: a \omterm{def:b1:graph-traversals:connected}{connected graph} has a \omterm{def:b1:graph-traversals:connected}{strongly connected} orientation if and only if it has no
\omterm{def:b1:graph-traversals:bridge}{bridge}. Brenmoor's road \omterm{def:b1:graph-traversals:graph}{graph} has no \omterm{def:b1:graph-traversals:bridge}{bridge}, and the orientation computed makes all 53 roads one-way with a single
\omterm{def:b1:graph-traversals:connected}{strongly connected component}, checked by Kosaraju's algorithm.
\end{solution}
